Exact statevector simulation of quantum circuits needs memory that doubles with every qubit. We distribute a complex128
statevector over the high bandwidth memory of one, two or four AMD Alveo U55C cards in one host and evaluate the design
on the physical cards. A single primitive, a $2\times2$ complex128 matrix applied to amplitude pairs under a per
operation slice enable, lets the cards execute every gate that needs no amplitudes of another card, including diagonal,
controlled and permutation gates on the partition qubits and wide multi controlled gates, while a compiler packs gates
into passes over memory and moves amplitudes only for the remaining gates. Each card runs \NumEnginesFlagship~engines, which use all \NumPcsUsedFlagship~memory
channels of the card. The system simulates \NumMaxWidthOneCard~qubits on one card, \NumMaxWidthTwoCards~on two and \NumMaxWidthFourCards~on four cards,
the largest widths its memory layout allows; all \NumCorrectnessInstances~checked instances agree with an independent reference to double precision rounding.
On the same cards the quantum Fourier transform at 20 qubits runs \NumSubmittedOverRevisedQftTwenty~times faster than
in a baseline design that computes the gates on partition qubits on the host, whose timing a microbenchmark model reproduces without fitted parameters. Ablations, timing
breakdowns, scaling, energy and host measurements show where the time goes: distribution adds capacity, not speed,
because exchanges between cards pass through host memory, and Qiskit Aer on the \NumHostCores~core host is faster for almost all
circuits measured.
